
\documentclass[11pt,a4paper]{moderncv}
\moderncvtheme[green]{classic}                
\usepackage[T1]{fontenc}
\usepackage[utf8x]{inputenc}
\usepackage[italian]{babel}
\usepackage[scale=0.8]{geometry}
\recomputelengths                             

\fancyfoot{} 

\renewcommand*{\namefont}{\fontsize{25}{28}\mdseries\upshape}

\firstname{Patricio Nicolas}
\familyname{Massaro Rocca}

\address{District of Columbia, USA}{}    
\mobile{+1-202-754-6052}                    
\email{patomassaro@gmail.com}                

\begin{document}
\maketitle
\vspace{-10mm}

\section{Work experience}

\cventry{2020 - 2024 }{Lead Machine Learning Engineer}{\href{https://www.bairesdev.com/}{Bairesdev}}{}{}{
\begin{itemize}
    \item Led the development of Machine Learning features using AWS, Databricks, Airflow, MLFlow, and SparkML, enhancing operational efficiency by generating 500M daily predictions and doubling marketing campaign conversion rates.
    \item Solved complex problems using Machine Learning to develop early detection algorithms, addressing data quality issues by reducing spam rates by 30\%.
    \item Developed visualizations for complex ML findings using Tableau to support strategic decision-making processes among top management.
    \item Designed and implemented a natural language processing deep learning model to classify inbound messages, reducing human workload by 37\%.
    \item Defined the data architecture to optimize ML models training, increasing performance metrics by 10\%.
\end{itemize} 
}

\cventry{2023-2024}{Thesis working student}{\href{https://ororatech.com/}{Ororatech}}{}{}{  
\begin{itemize}
    \item Applied Machine Learning using GANs in Pytorch across remote-sensing sectors, creating robust solutions that achieve a 20\% error reduction in land surface temperature estimation for wildfire management.
    \item Implemented a deep convolutional model for fine-scale georeferencing of thermal data products, improving accuracy in disaster monitoring applications.
    \item Published paper at IGARSS 2023 regarding multi-spectral remote sensing super-resolution.
\end{itemize} }

\cventry{2017 - 2020}{Business Analytics Engineer}{\href{https://www.tenaris.com/en}{Tenaris}}{}{}{
 \begin{itemize}
     \item Led the architectural design and workflow development for migrating the entire organization to a centralized cloud-based PowerBI application for 30,000 contributors.
     \item Managed the population of the data warehouse using Synapse and Data Factory and designed the semantic layer using Azure Analysis Services, reducing costs by 17\%.
     \item Defined platform adoption KPIs and action plans for their improvement.
 \end{itemize}
 }

\cventry{2017 - 2018}{Research Intern}{University of Buenos Aires}{}{}{ 
\begin{itemize}
     \item Two publications regarding the effects of piezoelectric devices in optoacoustic image reconstruction.
 \end{itemize}}

\cventry{2014 - 2017}{Systems Administrator}{Masstec S.R.L}{}{}{
    \begin{itemize}
        \item Administration of small businesses server, workstation, and networking systems.
    \end{itemize}
}

\cventry{2013- 2014}{Network Administrator}{Education Ministry, Buenos Aires}{}{}{
\begin{itemize}
    \item Engaged with educational users to refine ICT features and solutions, enhancing digital literacy for over 250 students through the 'Conectar Igualdad' program. 
\end{itemize}
}


\section{Education}
\cventry{2021-2023}{Elite MSc Data Science}{\href{https://www.m-datascience.mathematik-informatik-statistik.uni-muenchen.de/index.html}{Ludwig-Maximilians-Universität}}{Munich}{Germany}{Thesis: Blind Super Resolution Applied to Thermal Remote Sensing Data Products}
\cventry{2020-2021}{MSc Data Mining }{\href{http://datamining.dc.uba.ar/datamining/}{Universidad de Buenos Aires}}{}{Argentina}{}
\cventry{2010-2018}{Electronic Engineering}{\href{https://www.fi.uba.ar/}{Universidad de Buenos Aires}}{}{Argentina}{Thesis: Improvements in optoacoustic image reconstruction using deconvolution filters.} 

\section{Language}
\cvline{Spanish}{Native language.}
\cvline{English}{Bilingual proficiency.}
\cvline{German}{Professional working proficiency.}

\vspace{-3mm}

\end{document}
